\chapter{Reference}
\label{ch:reference}

\section{Files}

\begin{longtable}{@{}p{0.42\textwidth}p{0.52\textwidth}@{}}
\toprule
\textbf{File} & \textbf{Content}\\
\midrule
\endhead
\file{FX-Mechanics/fxmenet-server-plugins.xml} & The server's plugin cache.\\
\file{FX-Mechanics/fxmenet-server-settings.json} & The folders and the default-folders switch set with \cmd{fxmenet-top}.\\
\file{FX-Mechanics/FxmeNet.settings} & Client: the last server, the hover-help switch, the extra plugin folders (machine-wide).\\
\file{FX-Mechanics/fxmenet-client-plugins.xml} & Client: the cache of the local plugin scan.\\
\file{FX-Mechanics/fxmenet-client-events.log} & Client: the transport events log, when enabled.\\
\bottomrule
\end{longtable}

\noindent The \file{FX-Mechanics} folder is in the user's application-data
folder (on Linux, \file{~/.config}; on Windows,
\file{%APPDATA%}; on macOS, \file{~/Library/Application Support}). \todo{verify these locations on each system}

\section{Ports}

See chapter~\ref{ch:network}.

\section{Limits}

\begin{itemize}
  \item 256 automation parameters per \ui{FxmeNet} instance.
  \item 1 to 64 channels in and out of each \ui{FxmeNet} instance (stereo by
    default).
  \item Block sizes from 32 to 1024 samples; a network budget of 1 to 8 blocks.
  \item Plugins: VST3 effects on the server.
\end{itemize}

\section{More documentation}

The project's \file{doc} folder holds the technical documents:
\file{architecture.md} (how the code is built), \file{protocol.md} (the wire
protocol, for anyone writing another client or tool), \file{routing.md} (the
design of the routing) and \file{Tests/stress-tests.md} (the measurements of
chapter~\ref{ch:performance} in full).
